\begin{abstract}
\noindent In music streaming platforms, up to 35\% of user interaction logs correspond to immediate skips ($<30$ seconds), reflecting user dissatisfaction. Standard Maximum Likelihood Estimation (MLE) penalizes causal autoregressive transformers for failing to predict such negative interactions (the \textit{False Negative} dilemma), inadvertently amplifying disliked content. In this work, we investigate the alignment of a Causal GPT-2 sequential recommender (43.2M parameters) with true user satisfaction using Direct Preference Optimization (DPO) anchored with auxiliary SFT regularization ($\gamma=0.10$). Extensive evaluations on the industrial-scale Yandex Music dataset (\textbf{Yambda-50M}: 98,505 catalog tokens, 33.2M training interactions, 9,048 held-out test users) demonstrate that DPO not only outperforms 2-stage Critic Re-Ranking and Actor-Critic PPO, but completely circumvents the alignment tax (\textit{Zero Alignment Tax}). The aligned DPO policy achieves statistically significant next-item ranking gains (\textbf{+0.67\% HR@10}, $p < 0.0001$; McNemar's paired test $\chi^2 = 18.32, p = 1.9 \times 10^{-5}$), expands catalog coverage by \textbf{+1,027 unique tracks} from the long tail (+1.04\% of the catalog), reduces popularity Spearman correlation by \textbf{$-12.1\%$}, and yields a \textbf{+1.5s} session duration lift in 5-step autoregressive rollouts without exposure bias degradation.

\vspace{0.5em}
\noindent\textbf{Keywords:} Sequential Recommendation, Reinforcement Learning Alignment, Direct Preference Optimization, Causal GPT-2, Music Streaming, Skip Dilemma, Zero Alignment Tax, Catalog Diversity, Yambda.
\end{abstract}
